\subsection{A retained row changes the numerator chart}
The smallest selected endpoint uses a further exact algebraic adjustment,
not a weaker-frame theorem. Let $R_W=0$ be the retained frame residual of
the conservative source and replace the orientation numerator by
\[
 \Omega'=\Omega+R_W.
\]
All retained source zeros have $R_W=0$, so the joint zero-set argument is
unchanged. In the actual centered source coordinates this addition cancels
the quadratic $x^2$ term. A later monic substitution for $x$ can then have
degree four without making the selected pool row exceed degree eight.
The certificate checks the exact coefficient cancellation and restores the
retained row before invoking Theorem~\ref{thm:unweighted-joint}.
It is not a claim that $\Omega'$ is itself nonnegative away from those
retained zeros.

For the 113-variable endpoint the four base blocks are partitioned as
$(B)$ and $(D,q,q+1)$, with orientation fused into the first part. The
source has thirteen principal pairs and fifteen power requests, not the
fourteen-pair weighted reference. The 202-coordinate, 142-row joint seed
then undergoes 89 exact unit-pivot eliminations, leaving 113 coordinates
and 53 residuals of degree at most eight. Its sum of squares has exact
degree 16, support 1030202 and height $16250018885420533678080$.

The sparse 114-variable companion keeps orientation separate and partitions
the power blocks as $(B,D)$ and $(q,q+1)$. It has 54 residuals, exact degree
16, 127818 terms and the same height. This is a useful sparseness tradeoff,
not another use of the fused orientation theorem. The 138-variable companion
has 78 residuals, exact degree eight, 12099 terms and height $2^{50}$.
The 138-variable point uses direct monic elimination from the conservative
source, without a new pool, and is not credited to the joint theorem.
All three full collected polynomials are supplied; their supports are not
estimates inferred from circuit sizes.

\subsection{The finite substitution specification}
For clarity, a unit-pivot step at the current residual row has the form
\[
 Q_r=\epsilon x_v+f(\widehat{x_v}),\qquad\epsilon\in\{-1,1\},
\]
with no other occurrence of $x_v$ in $Q_r$. Substitute $x_v=-\epsilon f$
in every other residual, then delete row $r$ and coordinate $v$. This is
an integral zero-set bijection: insertion of that polynomial is its exact
inverse, and the deleted equation becomes an identity. Later restorations
compose in reverse order. Every step's resulting degree is checked after
full exact coefficient collection; a generic degree-propagation bound is
not used as an attained degree certificate.

The appendix gives the complete conservative source, literal principal-pair
and request ports, exact removed rows, fixed pool partitions, and each
ordered triple $(v,r,\epsilon)$. Starting with the stated seed, set
$\sigma_0$ to the identity and successively apply
\[
 \sigma_{j+1}=\sigma_j
 \bigl[x_v\mapsto-\epsilon\,\sigma_j(Q_r)|_{x_v=0}\bigr].
\]
Here the substitutions mean applying the current pivot replacement to every
current row; the recorded order removes any ambiguity about composition.
The final polynomial is the sum of squares of all retained substituted
rows, in the explicitly listed retained coordinate order. This is a finite
ordinary polynomial definition with no existential register or division.
The supplied full residual and coefficient arrays independently agree with
this specification.

The matched unfused control uses the same source and admissible monic pivot
strategy. It isolates the extra elimination enabled by the joint numerator
chart. The sparse 114-variable endpoint and the low-degree 138-variable
endpoint are separate controls with their own partitions and metrics; they
are not falsely reported as a simultaneous improvement in all objectives.
The unbounded joint theorem, exact substitutions and inherited source
calibration together establish~\eqref{eq:calibration} for every selected
endpoint on $\Z$.
